% 【本文件写什么】api-v0 契约层，叶子，只写语义不写实现
% - crate 路径：os/components/wateros-abi/abi-api/api-v0/src/
% - 列出并说明：pub trait、核心类型、错误枚举、常量
% - 每个公开 API 的前置条件、后置条件、线程/中断上下文限制
% - 与 Linux/用户态约定的对齐点与 deliberate 差异
% - v0 版本当前行为与后续可替换 extension 点
% 【事实来源】
% - 上述 crate 下全部 .rs 源文件
% - docs/exports/public-api/ 中对应符号表，以源码为准
% 【不写什么】
% - impl 内部数据结构、汇编、具体算法步骤

\paragraph{ErrNo}

\code{ErrNo(isize)}封装Linux errno正数值，与\code{asm-generic}数值一致。\code{raw()}返回正数errno；\code{user\_ret()}返回\code{-errno}供用户态\code{isize}承载。关联常量覆盖\code{EPERM}、\code{ENOENT}、\code{EINTR}、\code{EFAULT}、\code{EINVAL}、\code{ENOSYS}、\code{EAGAIN}、\code{ENOMEM}、\code{EACCES}、\code{EBADF}、\code{ECHILD}、\code{ENOSPC}、\code{ENOTTY}、\code{ENOTDIR}、\code{EISDIR}、\code{ENAMETOOLONG}、\code{ELOOP}、\code{ERANGE}等常用子集。\code{KernelResult<T>}为\code{Result<T, ErrNo>}别名。

\paragraph{UserRet 与 SyscallResult}

\code{UserRet(isize)}表示用户态可见的单一返回值：非负为成功，负值为\code{-errno}。\code{from\_success(usize)}、\code{from\_error(ErrNo)}、\code{from\_kernel\_result(SyscallResult)}完成编码。\code{SyscallResult = Result<usize, ErrNo>}是内核handler常用中间类型，由syscall层最终转为\code{UserRet}。

\paragraph{SyscallNumber 与 SyscallNumberTable}

\code{SyscallNumber(usize)}是调用号newtype，\code{new}/\code{raw}不校验内核是否实现。\code{SyscallNumberTable} trait以关联常量形式声明符号化调用号，按类别分组：文件与描述符、路径与xattr、进程与执行、内存、时间、凭证、信号、socket、多路复用等。trait只表达编号映射，不表示handler已实现；早期维护busybox/simple进程所需子集，随strace缺口逐步补齐。

\paragraph{SyscallArgs 与 SyscallPacket}

\code{SyscallArgs}为\code{repr(C)}的\code{[usize; MAX\_SYSCALL\_ARGS]}包装，当前 \code{MAX\_SYSCALL\_ARGS = 6}，顺序与ABI参数寄存器一致。\code{from\_regs}、\code{arg(idx)}、\code{as\_regs}提供构造与读写字段。\code{SyscallPacket}组合\code{SyscallNumber}与\code{SyscallArgs}，\code{new}不校验语义一致性。索引越界或未定义槽位语义属于调用方错误，由具体handler决定返回\code{EINVAL}或\code{EFAULT}。

\paragraph{与 Linux 对齐点}

错误码数值、成功非负/失败\code{-errno}约定、六个参数寄存器上限与Linux generic 64位用户态ABI对齐。deliberate差异：\code{SyscallNumberTable}未覆盖全部Linux调用；\code{SELECT}在generic64表上无独立nr，见 \code{impl-linux-generic64}。启用\code{impl-linux-generic64}时聚合层额外导出\code{ActiveSyscallNumberTable}类型别名。
